\input{LectureSlides/includes/metropolis_preamble}

% Title page info
\title[Dynamic Games]{Chapter 4: General Dynamic Games}
\author[ECON 5001]{ECON 5001\\Prof.Paul J. Healy\\Based on: \textit{Game Theory} by Giacomo Bonanno} \color{metrop}
\institute[OSU]{}
\date[]{\vfill {\tiny Updated \today\ at\ \DTMcurrenttime}}

% ---------------------------------------------------------------------
% Trees use xgames. New in this deck, beyond deck 05:
%   \informset (w1) -- (w2);       enclose nodes in one information set
%   \informset[dashed] / [curved=15]   alternative renderings
%   \subgame[inner sep=3ex] at (w2);   circle the subgame starting at w2
% Both must come AFTER the \branch commands that define those coordinates,
% and both are overlay aware, e.g. \subgame<3->.
% Put label=none on the second branch of an information set so the
% player's name isn't printed twice.
% ---------------------------------------------------------------------

\begin{document}

\frame{\maketitle}

\begin{frame}{What Chapter 3 Assumed}
Last time, every player knew the entire history when it was her turn.
\bsk
That ruled out almost everything interesting:
\begin{itemize}
    \item sealed bids
    \item a firm choosing capacity without seeing its rival's choice
    \item poker
    \item any move made \textbf{at the same time} as somebody else's
\end{itemize}
\vfill
In this lecture we drop that assumption, then rebuild backward induction so it
still works.
\vspace{-0.4in}
\end{frame}

% ==============================================================
% MOBLAB --- PLAY FIRST
% ==============================================================

% ---- MobLab run: coordination game -----------------------------------
% Game: "Matrix: Instructor Specified", 2x2, one shot, no chat. Payoffs:
%   Startup/Startup 3,3   Startup/Job 0,1   Job/Startup 1,0   Job/Job 1,1
% Label the actions "Startup" and "Corp job" if the interface allows,
% otherwise use A and B and map them aloud.
% Play BEFORE information sets are introduced. The same game returns twice:
% drawn as a tree with an information set (Section 4.1), and as the
% post-entry subgame in the SPE example (Section 4.4).
% ---------------------------------------------------------------------
\begin{frame}{Log In: Startup or Safe Job?}
\begin{itemize}
    \item You and a friend have an idea for an ambitious startup
    \item It only works if \textbf{both} of you quit and go all in
    \item The alternative is a cushy corporate job: safe, decent money
    \item Go all in alone and the startup fails. You get nothing.
\end{itemize}
\bsk
You two decide \textbf{at the same time}, with no chance to talk first.
\vfill
\centering
Pick one.
\vspace{-0.3in}
\end{frame}

\begin{frame}{What We Did}
\centering
\begin{matrixgame}[top={Startup, Corp job}, left={Startup, Corp job}, player names={You, Your friend}, br]
  \payoffmatrix{3,3 & 0,1 \\ 1,0 & 1,1}
\end{matrixgame}
\bsk
\raggedright
\begin{tabular}{lc}
Share who went all in: & \underline{\hspace{1in}} \\[0.25in]
Share who took the job: & \underline{\hspace{1in}} \\[0.25in]
\end{tabular}
\bsk
Two Nash equilibria: (Startup, Startup) and (Corp job, Corp job). One is
better for everyone. Getting there requires trusting your friend.
\vspace{-0.3in}
\end{frame}

% ==============================================================
% 4.1 IMPERFECT INFORMATION
% ==============================================================

\begin{frame}[standout]
Information Sets
\end{frame}

\begin{frame}{A Familiar Situation}
To discourage copying, I print two versions of the exam, white paper and pink,
and hand them out so neighbors never share one.
\bsk
Or do I? Writing genuinely different versions is extra work for me. I might be
bluffing: same questions, different paper.
\bsk
A student who didn't study is sitting next to someone who did.
\begin{itemize}
    \item Turn in a blank exam and, given earlier grades, he gets a \textbf{C}
    \item Copy, and if the versions are identical he gets an \textbf{A}
    \item Copy, and if the versions differ he is caught and gets an \textbf{F}
\end{itemize}
\vfill
The moves are sequential. I decide first, but he doesn't \textit{see} my decision.
\vspace{-0.3in}
\end{frame}

\begin{frame}{Drawing What He Doesn't Know}
\centering
\begin{gametree}[player names={Professor, Student}, scale=0.85]
  \branch[player=1, root=w, label=90] from (0,0) to[v=1.3, h=6.2]
    {Identical versions; Different versions};
  \branch[player=2, label=none] from (w1) to[v=1.5, h=2.6]
    {Copy(\makecell{Student\\ gets A}); Blank(\makecell{Student\\ gets C})};
  \branch from (w2) to
    {Copy(\makecell{Caught:\\ gets F}); Blank(\makecell{Student\\ gets C})};
  \informset (w1) -- (w2);
\end{gametree}
\bsk
\raggedright
The enclosure is an \textbf{information set}: the Student knows he is at one of
those two nodes, but not which one.
\vspace{-0.2in}
\end{frame}

\begin{frame}{Definition: Extensive-Form Game-Frame}
Everything from Chapter 3, plus one item:
\bsk
\begin{itemize}
    \item a finite rooted directed tree
    \item players $I=\{1,\ldots,n\}$, one assigned to every decision node
    \item actions on edges, no two edges out of a node sharing an action
    \item outcomes assigned to terminal nodes
\end{itemize}
\bsk
$\star$ \ For every player $i$, a \textbf{partition} $\mathcal{D}_i$ of her decision
nodes $D_i$. Each cell of the partition is an \textbf{information set}.
\bsk
We write $H$ for a single information set, so $H\in\mathcal{D}_i$.
\bsk
\centering
\textbf{Restriction:} any two nodes in the same information set must offer
\textit{exactly the same actions}.
\vfill
\raggedright
Why must that hold? \vspace{0.3in}
\end{frame}

\begin{frame}{A Note on Notation}
Three different things, easily confused:
\bsk
\begin{itemize}
    \item $D_i$ \ is the set of \textbf{all} decision nodes belonging to player $i$
      \vspace{0.15in}
    \item $\mathcal{D}_i$ \ is the \textbf{partition} of $D_i$, a collection of sets
      \vspace{0.15in}
    \item $H\in\mathcal{D}_i$ \ is \textbf{one} information set, a single cell
\end{itemize}
\vfill
{\footnotesize Bonanno writes $D$ and $D'$ where we write $H$ and $H'$. Watch out:
in Chapter 3 he also used $D$ for the set of \textit{all} decision nodes in the
tree, so the same letter does two jobs. We'll stick with $H$ for information sets.}
\vspace{-0.2in}
\end{frame}

\begin{frame}{Perfect Information Is the Special Case}
\vfill
\centering
\large
Every information set is a singleton\\
\bsk
$\Longleftrightarrow$\\
\bsk
Perfect information
\vfill
\normalsize
\raggedright
So Chapter 3 was not a different theory. It was this one, with the partition
as fine as it can be.
\bsk
Convention from here: we draw an enclosure \textbf{only} around information sets
with two or more nodes. A lone node is its own information set.
\vspace{-0.4in}
\end{frame}

\begin{frame}{Perfect Recall: First Condition}
Almost all of game theory assumes players don't forget.
\bsk
\begin{block}{Condition 1}
No node in an information set is a \textbf{predecessor} of another node in it.
\end{block}
\bsk
Below, $H=\{r,u\}$ violates it: $r$ comes before $u$, so Player 1 cannot tell
whether she has already moved.
\vfill
\centering
\begin{gametree}[labels=short, scale=0.78]
  \branch[player=1, root=w, label=90, label text=$r$] from (0,0) to[v=1.3, h=4.0] {a; b(\ )};
  \branch[label=180, label text=$u$] from (w1) to[v=1.3, h=1.6] {a(\ ); b(\ )};
  \informset[curved=25, label=none] (w) -- (w1);
\end{gametree}
\vspace{-0.15in}
\end{frame}

\begin{frame}{Perfect Recall: Second Condition}
\begin{block}{Condition 2}
If $x\in H$ leads to $y\in H'$ by action $a$, then \textbf{every} node in $H'$
has a predecessor in $H$ reached by that same $a$.
\end{block}
\bsk
Below, $H'$ is Player 2's lower information set. Its right node is reached via
$H$, but its left node never passes through $H$ at all. Player 2 cannot tell
whether this is her first move or her second.
\vfill
\centering
\begin{gametree}[labels=short, scale=0.7]
  \branch[player=1, root=w, label=90] from (0,0) to[v=1.4, h=6.4] {a; b};
  \branch[player=2, label=0, label text=$H$] from (w2) to[v=1.4, h=2.6] {c; d(\ )};
  \branch[label=none] from (w1) to[v=1.4, h=1.6] {e(\ ); f(\ )};
  \branch[label=none] from (w21) to {e(\ ); f(\ )};
  \informset[curved=18] (w1) -- (w21) {$H'$};
\end{gametree}
\vspace{-0.15in}
\end{frame}

\begin{frame}{The Same Game, Drawn as a Tree}
\centering
\only<1>{\begin{gametree}[player names={You, Your friend}, scale=0.85]
  \branch[player=1, root=w, label=90] from (0,0) to[v=1.3, h=5.0] {Startup; Corp job};
  \branch[player=2, label=none] from (w1) to[v=1.4, h=2.0] {Startup(3,3); Corp job(0,1)};
  \branch from (w2) to {Startup(1,0); Corp job(1,1)};
  \informset (w1) -- (w2);
\end{gametree}}%
\only<2>{\begin{gametree}[player names={You, Your friend}, scale=0.85]
  \branch[player=2, root=w, label=90] from (0,0) to[v=1.3, h=5.0] {Startup; Corp job};
  \branch[player=1, label=none] from (w1) to[v=1.4, h=2.0] {Startup(3,3); Corp job(1,0)};
  \branch from (w2) to {Startup(0,1); Corp job(1,1)};
  \informset (w1) -- (w2);
\end{gametree}}
\bsk
\raggedright
\onslide<1->{One of you is drawn as moving first, but the other's information set
means she sees nothing. So this is \textbf{the same game} as the matrix.}
\bsk
\onslide<2>{Swap who moves first and nothing changes. Simultaneity is not about
time. It is about \textbf{what you know when you choose}.}
\vspace{-0.2in}
\end{frame}

% ==============================================================
% 4.2 STRATEGIES
% ==============================================================

\begin{frame}[standout]
Strategies, Again
\end{frame}

\begin{frame}{One Choice per Information Set}
\begin{block}{Definition}
A \textbf{strategy} for a player in an extensive-form game is a list of choices,
one for every \textbf{information set} of that player.
\end{block}
\bsk
Chapter 3 said ``one for every decision node.'' That was the same rule, because
every node was its own information set.
\vfill
The change matters: a player can no longer make a plan that depends on
something she cannot observe.
\vspace{-0.4in}
\end{frame}

\begin{frame}{A Plan You Are Not Allowed to Make}
Two candidates, Yvonne and Fran, interview for one position. The employer says:
``don't call me, I'll call you.'' If the first candidate turns him down, he
immediately calls the other, never saying it's a second offer.
\bsk
Yvonne would love to plan:
\textit{If he calls me first, I say Yes. If I'm the recycled offer, I say No.}
\bsk
\centering
\begin{gametree}[player names={Employer, Yvonne, Fran}, scale=0.58]
  \branch[player=1, root=w, label=90] from (0,0) to[v=1.6, h=8.4]
    {Call Yvonne; Call Fran};
  \branch[player=2, label=none] from (w1) to[v=1.7, h=2.2] {Yes(\ ); No(\ )};
  \branch[player=3, label=0] from (w2) to[v=1.7, h=3.0] {Yes(\ ); No};
  \branch[player=2, label=none] from (w22) to[v=1.7, h=2.2] {Yes(\ ); No(\ )};
  \informset[player=2, curved=16] (w1) -- (w22);
\end{gametree}
\bsk
\raggedright
Must choose the \textbf{same action at every node} of an information set.
\vspace{-0.2in}
\end{frame}

\begin{frame}{From Tree to Table, Again}
Every strategy profile still determines a path, hence an outcome, hence payoffs.
So every extensive-form game still has a \textbf{strategic form}.
\bsk
And we could just find its Nash equilibria and stop.
\vfill
But Chapter 3 taught us not to. Some Nash equilibria of the strategic form rest
on threats nobody would carry out.
\bsk
We fixed that with backward induction, which needs a \textit{last} decision node.
\vfill
\centering
\large
What is the analogue when players move blind?
\vspace{-0.3in}
\end{frame}

% ==============================================================
% 4.3 SUBGAMES
% ==============================================================

\begin{frame}[standout]
Subgames
\end{frame}

\begin{frame}{The Idea}
A \textbf{subgame} is a piece of the tree that could stand alone as a game.
\bsk
Informal procedure:
\begin{enumerate}
    \item Pick a decision node $x$ whose information set is \textbf{just} $\{x\}$
    \item Draw an oval around $x$ and everything that follows it
\end{enumerate}
\bsk
If that oval \textbf{cuts through} an information set, the oval is illegal.
So only call it a subgame if it cuts \textbf{no} information set.
\bsk
``Cutting'' means some information set has one node inside your oval and another
node outside it.
\vfill
Every game is a subgame of itself. Anything smaller is a \textbf{proper} subgame.
\vspace{-0.3in}
\end{frame}

\begin{frame}{Why Cutting Is Fatal}
\centering
\begin{gametree}[labels=short, scale=0.85]
  \branch[player=1, root=w, label=90] from (0,0) to[v=1.3, h=5.6] {L; R};
  \branch[player=2, label=135] from (w1) to[v=1.5, h=1.8] {a(2,2); b(0,0)};
  \branch[player=1, label=45] from (w2) to[v=1.5, h=2.4] {U; D(2,0)};
  \branch[player=2, label=none] from (w21) to[v=1.5, h=1.6] {a(1,3); b(3,1)};
  \informset[curved=15] (w1) -- (w21);
\end{gametree}
\bsk
\raggedright
Player 1's node after $R$ looks like a fine place to start. But the oval around
it would contain \textit{one} of Player 2's two nodes.
\bsk
Player 2 can't tell those nodes apart, so you cannot hand her the right-hand
piece and call it a game. Half her uncertainty is outside it.
\vspace{-0.2in}
\end{frame}

\begin{frame}{Definition: Subgame}
Let $S(x)$ be the set containing $x$ and all its successors.
\bsk
$S(x)$ is a \textbf{subgame starting at $x$} if:
\bsk
\begin{enumerate}
    \item the information set containing $x$ is the singleton $\{x\}$, and
    \bsk
    \item for every information set $H$ in the game, \ if $H\cap S(x)\ne\emptyset$
      \ then \ $H\subset S(x)$
\end{enumerate}
\bsk
It is a \textbf{proper} subgame if $x$ is not the root.
\vfill
A subgame is \textbf{minimal} if it contains no other subgame strictly inside it.
Those are the ones we solve first.
\vspace{-0.3in}
\end{frame}

% ==============================================================
% 4.4 SUBGAME-PERFECT EQUILIBRIUM
% ==============================================================

\begin{frame}[standout]
Subgame-Perfect Equilibrium
\end{frame}

\begin{frame}{Definition}
Let $s$ be a strategy profile for the whole game, and $G$ a subgame. Write
$s|_G$ for the \textbf{restriction} of $s$ to $G$: the part of $s$ that prescribes
choices at $G$'s information sets, and only those.
\bsk
\begin{block}{Definition}
$s$ is a \textbf{subgame-perfect equilibrium} if, for \textbf{every} subgame $G$,
the restriction $s|_G$ is a Nash equilibrium of $G$.
\end{block}
\bsk
The whole game is one of its own subgames, so every subgame-perfect equilibrium
is a Nash equilibrium.
\vfill
Not every Nash equilibrium is subgame perfect. That's the point.
\vspace{-0.4in}
\end{frame}

\begin{frame}{IKEA and Wayfair, Again}
IKEA is the established furniture store at Polaris, earning \$10 million a year
in profit. Wayfair decides whether to open across I-71.
\bsk
If Wayfair opens, the two stores \textbf{simultaneously} decide whether to fund a
joint campaign advertising Polaris as a furniture destination.
\begin{itemize}
    \item Both fund it: shoppers drive up from all over. Both gain.
    \item Neither funds it: nothing changes.
    \item One funds it alone: too small to work, and the money is gone.
\end{itemize}
\bsk
\vfill
That second stage is the startup game you played earlier.
\vspace{-0.3in}
\end{frame}

\begin{frame}{The Whole Game}
\centering
\begin{gametree}[player names={Wayfair, IKEA}, scale=0.72, subgame={depth=9mm}]
  \branch[player=1, root=w, label=90] from (0,0) to[v=1.6, h=7.6] {Out(2,10); In};
  \branch[player=1, label=0] from (w2) to[v=1.6, h=4.0] {Fund; Skip};
  \branch[player=2, label=none] from (w21) to[v=1.6, h=1.5]
    {Fund(3,3); Skip(0,1)};
  \branch from (w22) to {Fund(1,0); Skip(1,1)};
  \informset (w21) -- (w22);
  \subgame[inner sep=4.6ex] at (w2);
\end{gametree}
\bsk
\raggedright
{\footnotesize Annual profit in \$millions. First number Wayfair, second IKEA.}
\ The circle marks the only proper subgame.
\vspace{-0.2in}
\end{frame}

\begin{frame}{The Algorithm}
\begin{enumerate}
    \item Find a \textbf{minimal} subgame.
    \bsk
    \item Solve it. Find its Nash equilibria, pick one, and note the strategies.
    \bsk
    \item \textbf{Mark} the subgame's starting node with that equilibrium's payoff
      vector, exactly as in backward induction, then delete everything below it.
      You now have a smaller game.
    \bsk
    \item Repeat until no proper subgames remain. Solve what's left. Patch together
      everything you noted along the way.
\end{enumerate}
\bsk
{\footnotesize Same marking idea as Chapter 3. The only change: a minimal subgame may
be a whole simultaneous game rather than a single node, so we mark it with a Nash
equilibrium payoff instead of one player's best choice.}
\vfill
If a subgame had several Nash equilibria, run the procedure again with a different
one for a different subgame-perfect equilibrium.
\bsk
{\footnotesize If some subgame has \textit{no} Nash equilibrium, neither does the game.}
\vspace{-0.3in}
\end{frame}

\begin{frame}{Step 1: Solve the Subgame}
\centering
\begin{matrixgame}[top={Fund, Skip}, left={Fund, Skip}, player names={Wayfair, IKEA}, br]
  \payoffmatrix{3,3 & 0,1 \\ 1,0 & 1,1}
\end{matrixgame}
\bsk
\raggedright
Two Nash equilibria: \ (Fund, Fund) with payoffs $(3,3)$, \ and
(Skip, Skip) with payoffs $(1,1)$.
\vfill
So we have to run the algorithm \textbf{twice}.
\vspace{-0.4in}
\end{frame}

\begin{frame}{Step 2: Replace and Solve}
Using the \textbf{(Fund, Fund)} equilibrium.
\bsk
\centering
\begin{gametree}[player names={Wayfair, IKEA}, scale=0.72, subgame={depth=9mm}]
  \branch[player=1, root=w, label=90] from (0,0) to[v=1.6, h=7.6] {Out(2,10); In};
  \branch[player=1, label=0] from (w2) to[v=1.6, h=4.0] {Fund; Skip};
  \branch[player=2, label=none] from (w21) to[v=1.6, h=1.5]
    {Fund(3,3); Skip(0,1)};
  \branch from (w22) to {Fund(1,0); Skip(1,1)};
  \informset (w21) -- (w22);
  \subgame[inner sep=4.6ex] at (w2);
\end{gametree}
\vspace{-0.15in}
\end{frame}

\begin{frame}{Step 2: Replace and Solve}
Using the \textbf{(Skip, Skip)} equilibrium.
\bsk
\centering
\begin{gametree}[player names={Wayfair, IKEA}, scale=0.72, subgame={depth=9mm}]
  \branch[player=1, root=w, label=90] from (0,0) to[v=1.6, h=7.6] {Out(2,10); In};
  \branch[player=1, label=0] from (w2) to[v=1.6, h=4.0] {Fund; Skip};
  \branch[player=2, label=none] from (w21) to[v=1.6, h=1.5]
    {Fund(3,3); Skip(0,1)};
  \branch from (w22) to {Fund(1,0); Skip(1,1)};
  \informset (w21) -- (w22);
  \subgame[inner sep=4.6ex] at (w2);
\end{gametree}
\vspace{-0.15in}
\end{frame}

\begin{frame}{Two Subgame-Perfect Equilibria}
$$\big((\text{In},\text{Fund}),\ \text{Fund}\big)
\qquad\text{and}\qquad
\big((\text{Out},\text{Skip}),\ \text{Skip}\big)$$
\vfill
In the second one Wayfair never opens, \textit{because} it expects no campaign.
\bsk
Nothing about the market changed. The pessimism is self-confirming.
\vfill
\begin{itemize}
    \item Both firms would prefer the first equilibrium
    \item Neither can get there by acting alone
    \item Subgame perfection does \textbf{not} pick between them
\end{itemize}
\bsk
Refinements narrow the set of predictions. Not always to one.
\vspace{-0.4in}
\end{frame}

\begin{frame}{The Strategic Form}
Wayfair's strategy is a pair: what to do at the root, and what to do in the
campaign stage.
\bsk
\centering
{\small
\only<1>{\begin{matrixgame}[top={Fund, Skip}, left={(Out\,Fund), (Out\,Skip), (In\,Fund), (In\,Skip)}, player names={Wayfair, IKEA}]
  \payoffmatrix{2,10 & 2,10 \\ 2,10 & 2,10 \\ 3,3 & 0,1 \\ 1,0 & 1,1}
\end{matrixgame}}%
\only<2>{\begin{matrixgame}[top={Fund, Skip}, left={(Out\,Fund), (Out\,Skip), (In\,Fund), (In\,Skip)}, player names={Wayfair, IKEA}, br]
  \payoffmatrix{2,10 & 2,10 \\ 2,10 & 2,10 \\ 3,3 & 0,1 \\ 1,0 & 1,1}
\end{matrixgame}}}
\bsk
\raggedright
\onslide<2>{Three cells have both payoffs underlined. Three Nash equilibria.}
\vspace{-0.2in}
\end{frame}

\begin{frame}{The Nash Equilibrium We Threw Out}
The strategic form has \textbf{three} Nash equilibria. Wayfair's strategy is a pair:
what to do at the root, and what to do in the campaign stage.
\bsk
\centering
\begin{tabular}{lll}
\textbf{Profile} & \textbf{Subgame play} & \textbf{Subgame perfect?} \\
\hline
$\big((\text{In},\text{Fund}),\text{Fund}\big)$ & (Fund, Fund) & yes \\
$\big((\text{Out},\text{Skip}),\text{Skip}\big)$ & (Skip, Skip) & yes \\
$\big((\text{Out},\text{Fund}),\text{Skip}\big)$ & (Fund, Skip) & \textbf{no} \\
\end{tabular}
\bsk
\raggedright
\vfill
The third is a Nash equilibrium of the whole game. Wayfair staying out is a best
reply if it believes IKEA will skip.
\bsk
But its plan for the subgame is incoherent: if entry happened, funding alone
earns 0, and it would rather skip. \vspace{0.3in}
\end{frame}

\begin{frame}{Three Things Worth Remembering}
\begin{itemize}
    \item In a \textbf{perfect-information} game, subgame-perfect equilibrium is
      exactly backward induction. Every decision node starts a subgame, so the
      algorithm walks the tree from the leaves. \vspace{0.25in}
    \item If a game has \textbf{no proper subgames}, every Nash equilibrium is
      subgame perfect. The refinement has nothing to bite on. \vspace{0.25in}
    \item Otherwise subgame perfection is a strict refinement: fewer predictions,
      all of them Nash. \vspace{0.25in}
\end{itemize}
\vfill
{\footnotesize Bonanno's Figures 4.8 and 4.9 work through larger three-player trees
with several subgames. Worth doing before the problem set.}
\vspace{-0.3in}
\end{frame}

% ==============================================================
% 4.5 CHANCE MOVES
% ==============================================================

\begin{frame}[standout]
Chance Moves
\end{frame}

\begin{frame}{When Nobody Chose It}
Sometimes what happens isn't anyone's decision. A card is dealt, a part fails,
a customer walks in.
\bsk
We model this with a fictitious player: \textbf{Nature} (or \textbf{Chance}).
\bsk
\begin{itemize}
    \item Nature moves at some nodes
    \item Each of Nature's actions carries a \textbf{probability}
    \item Nature has no payoffs and no objectives, so it isn't really a player
\end{itemize}
\vfill
Everything else (information sets, strategies, subgames) is unchanged.
\vspace{-0.4in}
\end{frame}

\begin{frame}{Three Cards}
Three cards: one black, two red. Shuffled, face down.
\bsk
\textbf{Nature moves first}, turning up the top card: black with probability
$\tfrac{1}{3}$, red with probability $\tfrac{2}{3}$.
\bsk
Then the players move:
\begin{itemize}
    \item Adele looks at the top card without showing Ben
    \item She tells Ben ``black'' or ``red'', truthfully or not
    \item Ben then guesses the true color
    \item Correct guess: Adele pays Ben \$9. Wrong: Ben pays Adele \$9.
\end{itemize}
\vfill
Ben has two information sets, one for each thing Adele might say, and in each
he is unsure what the card actually is.
\vspace{-0.3in}
\end{frame}

\begin{frame}{The Card Game}
\centering
\begin{gametree}[player names={Adele, Ben}, scale=0.58]
  \branch[player=0, root=w, label=90] from (0,0) to[v=1.5, h=8.6]
    {$\tfrac13$ Black; $\tfrac23$ Red};
  \branch[player=1, label=135] from (w1) to[v=1.5, h=4.2] {``B''; ``R''};
  \branch[label=45] from (w2) to {``R''; ``B''};
  \branch[player=2, label=none] from (w11) to[v=1.9, h=1.5] {B(-9,9); R(9,-9)};
  \branch[label=none] from (w12) to {B(-9,9); R(9,-9)};
  \branch[label=none] from (w21) to {B(9,-9); R(-9,9)};
  \branch[label=none] from (w22) to {B(9,-9); R(-9,9)};
  \informset[dashed] (w12) -- (w21);
  \informset[curved=52] (w11) -- (w22);
\end{gametree}
\bsk
\raggedright
{\footnotesize Nature moves first. Adele's label is what she \textit{says}; Ben's is
what he \textit{guesses}. First payoff Adele, second Ben.}
\bsk
Ben has two information sets, one for each thing Adele might say. In each, he
knows what he heard but not what the card is.
\vspace{-0.2in}
\end{frame}

\begin{frame}{A New Difficulty}
Take the profile: Adele always tells the truth, Ben always guesses black.
\bsk
What's the outcome? It depends on the card, so the outcome is itself a gamble:
\bsk
\centering
\begin{tabular}{|l|c|c|}
\hline
Outcome & Adele pays \$9 & Ben pays \$9 \\
\hline
Probability & $\tfrac{1}{3}$ & $\tfrac{2}{3}$ \\
\hline
\end{tabular}
\bsk
\raggedright
\vfill
We call this a \textbf{lottery}. To build a strategic form we need to know how
players rank \textit{lotteries}, not just outcomes.
\bsk
Chapter 5 answers that properly. For now, one shortcut.
\vspace{-0.3in}
\end{frame}

\begin{frame}{Expected Value}
\begin{block}{Definition}
The \textbf{expected value} of the money lottery
$\begin{pmatrix} \$x_1 & \$x_2 & \cdots & \$x_n \\ p_1 & p_2 & \cdots & p_n\end{pmatrix}$
is
$$x_1p_1 + x_2p_2 + \cdots + x_np_n$$
\end{block}
\bsk
\centering
$\begin{pmatrix} \$5 & \$15 & \$25 \\[2pt] \tfrac{1}{5} & \tfrac{2}{5} & \tfrac{2}{5}\end{pmatrix}$
\ has expected value \ \underline{\hspace{0.9in}}
\bsk
$\begin{pmatrix} -\$9 & \$9 \\[2pt] \tfrac{1}{3} & \tfrac{2}{3}\end{pmatrix}$
\ has expected value \ \underline{\hspace{0.9in}}
\vfill
\end{frame}

\begin{frame}{Risk Neutrality}
\begin{block}{Definition}
A player is \textbf{risk neutral} if she considers a money lottery to be
\textit{just as good as} its expected value.
\end{block}
\bsk
So a risk-neutral player ranks lotteries by their expected value, and we can use
that number as her payoff.
\bsk
\centering
$L_1=\begin{pmatrix} \$5 & \$15 & \$25 \\[2pt] \tfrac{1}{5} & \tfrac{2}{5} & \tfrac{2}{5}\end{pmatrix}$
\hspace{0.2in}
$L_2=\begin{pmatrix} \$16 \\[2pt] 1\end{pmatrix}$
\hspace{0.2in}
$L_3=\begin{pmatrix} \$0 & \$32 & \$48 \\[2pt] \tfrac{5}{8} & \tfrac{1}{8} & \tfrac{1}{4}\end{pmatrix}$
\bsk
\raggedright
A risk-neutral player ranks these: \underline{\hspace{2in}}
\vfill
{\footnotesize This is a convenience, not a claim about people. Risk neutrality
lets us finish Chapter 4 without the machinery of Chapter 5.}
\vspace{-0.2in}
\end{frame}

\begin{frame}{Why the Assumption Is Doing Work}
$L_2$ and $L_3$ have the same expected value, \$16.
\bsk
\begin{itemize}
    \item $L_2$ pays \$16 for certain
    \item $L_3$ pays nothing five times out of eight, and sometimes \$48
\end{itemize}
\bsk
A risk-neutral player is \textbf{indifferent} between them.
\vfill
Most people are not. Ask yourself which you'd take, and notice that your answer
is information the expected value throws away.
\bsk
Recovering it is exactly what Chapter 5 is for.
\vspace{-0.4in}
\end{frame}

% ==============================================================
% WORKED EXAMPLE
% ==============================================================

\begin{frame}[standout]
Let's Work One Out
\end{frame}

\begin{frame}{A Three-Player Dynamic Game}
\centering
\begin{gametree}[player names={Ann, Ben, Cal}, scale=0.66]
  \branch[player=1, root=w, label=90] from (0,0) to[v=1.5, h=8.0] {Out(3,3,3); In};
  \branch[player=2, label=0] from (w2) to[v=1.5, h=4.2] {L; R};
  \branch[player=3, label=none] from (w21) to[v=1.5, h=1.8] {u(4,1,1); d(0,0,4)};
  \branch[label=none] from (w22) to {u(1,4,0); d};
  \branch[player=1, label=0] from (w222) to[v=1.5, h=1.8] {p(2,2,2); q(5,1,1)};
  \informset[player=3] (w21) -- (w22);
\end{gametree}
\vspace{-0.15in}
\end{frame}

\begin{frame}{A Three-Player Dynamic Game}
Its strategic form. Ann's strategy is a pair: the root, then her later node.
\bsk
\centering
{\scriptsize
\begin{matrixgame}[top={L, R}, left={(Out\,p), (Out\,q), (In\,p), (In\,q)}, bottom={Cal: u, Cal: d}, player names={Ann, Ben, Cal}]
  \payoffmatrix{3,3,3 & 3,3,3 \\ 3,3,3 & 3,3,3 \\ 4,1,1 & 1,4,0 \\ 4,1,1 & 1,4,0}
  \payoffmatrix[pos={4.6,0}]{3,3,3 & 3,3,3 \\ 3,3,3 & 3,3,3 \\ 0,0,4 & 2,2,2 \\ 0,0,4 & 5,1,1}
\end{matrixgame}}
\vspace{-0.15in}
\end{frame}

\begin{frame}{A Three-Player Dynamic Game}
\begin{enumerate}
    \item List every subgame. Why is there none starting at Cal's left node? \vspace{0.3in}
    \item How many strategies does each player have? \vspace{0.3in}
    \item Run the subgame-perfect equilibrium algorithm. Solve the smallest
      subgame first. \vspace{0.35in}
    \item Find a Nash equilibrium of the whole game that is \textbf{not}
      subgame perfect. \vspace{0.3in}
\end{enumerate}
\end{frame}

\begin{frame}{What Have We Learned, Charlie Brown?}
\begin{itemize}
    \item An \textbf{information set} collects nodes a player cannot tell apart;
      perfect information is the case where all are singletons
    \item A \textbf{strategy} gives one choice per information set, not per node
    \item A \textbf{subgame} starts at a singleton information set and cuts no
      information set
    \item \textbf{Subgame-perfect equilibrium}: Nash in every subgame. Solve minimal
      subgames first, replace, repeat
    \item \textbf{Nature} handles chance; risk neutrality ranks lotteries by mean
\end{itemize}
\vfill
Next: what if players aren't risk neutral?
\vspace{-0.4in}
\end{frame}

\begin{frame}{Reading \& Practice}
\begin{itemize}
    \item Bonanno, \textit{Game Theory}, Chapter 4
    \item Exercises: \S4.6.1 (imperfect information), \S4.6.2 (strategies),\\
      \S4.6.3 (subgames), \S4.6.4 (subgame-perfect equilibrium),\\
      \S4.6.5 (chance moves)
    \item Exercise 4.2 is the entry game with a simultaneous second stage, close to
      the IKEA example, with more to it
    \item Solutions in \S4.7
\end{itemize}
\vfill
\end{frame}

\end{document}
